%% sections/design.tex
\section{Experimental Design}
\label{sec:design}

\subsection{Session structure: four consecutive rounds}

The unit of study in~\citet{dufwenberg2005bubbles} is a
\emph{session} of four consecutive ten-period markets played by the
same agent population. The paper insists on the terminology explicitly:
\begin{quote}
``A session involved four consecutive markets. In the following, we
shall talk in terms of four different rounds. Note the distinction
between rounds and periods; a round (being a market) consists of
ten periods.''~\citep[p.~1733]{dufwenberg2005bubbles}
\end{quote}
The session wrapper is the load-bearing element of the original
design: in the human-subject study, bubbles inflate in round~1,
attenuate in rounds~2--3, and have effectively died by round~4 even
though the population is identical---experience, not composition,
kills the bubble. The testbed inherits this wrapper because it is
also the channel through which Plan~I's experience-weighted trust
runs: every surviving utility agent's \texttt{roundsPlayed} counter
is incremented at the end of each non-final round, which feeds
equation~\eqref{eq:plan1_body} and ramps the Plan~I belief-blend
weight from $0.60$ to $0.90$ across the session. The simulator
therefore wraps the per-round market in a session loop of $R=4$
rounds. At every round boundary the order book is cleared, each
agent's cash and inventory are rewound to its sampled spec, and the
period counter resets to one. Learned state---trust matrices, the
experience counter, the candidate action cache---is intentionally
preserved across the boundary so that the cross-round learning
channel that drives Plan~I's weight ramp is the only source of
persistent change between consecutive markets. Both
$R$ and $T$ are surfaced read-only in the simulator's ``Paper
constants'' panel and are not user tunable.

\subsection{The asset and the fundamental}

The asset is a finite-lived claim on an iid per-period dividend
that pays either zero or twenty experimental cents with equal
probability. Letting $d$ denote the period-$t$ dividend draw,
\begin{equation}
d\sim\text{Uniform}\{0,20\},\qquad \E[d]=10,\qquad T=10.
\label{eq:dividend}
\end{equation}
The asset lives for $T=10$ periods within each round and pays
nothing thereafter, so its risk-neutral fundamental value at the
start of period~$t$ of any round is
\begin{equation}
\FV_t \;=\; \E[d]\cdot(T-t+1) \;=\; 10\cdot(T-t+1),
\label{eq:fv}
\end{equation}
that is, a linearly declining staircase from $\FV_1=100$ to
$\FV_T=10$ with a terminal value of zero after the final dividend
pays. The fundamental resets to $\FV_1=100$ at the start of every
new round, so the price chart traces a four-tooth saw-tooth across
the full session. Equation~\eqref{eq:dividend} and
equation~\eqref{eq:fv} reproduce the dividend process and the
fundamental schedule of~\citet{dufwenberg2005bubbles} exactly. The
session length $R$, the per-round asset life $T$, and the dividend
mean $\E[d]$ are all surfaced read-only in the simulator's ``Paper
constants'' panel so that every comparison across plans holds the
market structure fixed.

\subsection{Continuous double auction with period discretisation}

The simulator discretises each period into $K=18$ decision rounds,
so a single round comprises $T\cdot K=180$ ticks and a full
four-round session comprises $R\cdot T\cdot K=720$ ticks. Every
tick, each agent posts at most one order, routed to a price--time
priority order book; a crossing pair matches at the resting price.
At the end of each period the outstanding quotes are withdrawn, the
dividend is drawn from~\eqref{eq:dividend} and credited to each
holder, and the belief-update plan (Plan~I, II, or~III) runs once
to write next period's $V_i^{\text{prior}}$. At the end of period~$T$
of every non-final round the session loop also rewinds endowments,
clears the book, increments every surviving agent's
\texttt{roundsPlayed} counter, and resets the period counter to one
before the next round begins. The $K=18$ discretisation is the
simulator's own approximation of the two-minute continuous z-Tree
window used in the paper, is pinned in the ``Simulator constants''
panel, and is not user tunable.

\subsection{Population and endowment sampling}
\label{sec:sampling}

The population is a mix of fundamentalists~(F), trend
followers~(T), zero-intelligence randoms~(R), and utility
agents~(U), dialled from the Experiment Settings panel. A utility
slot is always present because Plans~I, II, and~III act only on~U
agents---the background slots (F, T, R) are there to provide a
rational anchor, a momentum amplifier, and a liquidity floor, and
play the same role under every plan.

Before the engine starts, \texttt{sampleAgents(mix, rng)} draws a
flat list of agent specs, each carrying a numeric id, an agent
type label, a personal name drawn without replacement from a
built-in pool, a cash endowment drawn uniformly from
$[800,1200]$~cents, and an inventory drawn uniformly from
$\{2,3,4\}$. Utility agents additionally carry a persistent bias
mode~$b_i$ and risk-preference type~$U_i$, assigned from a
user-controlled split over optimistic/unbiased/pessimistic and
loving/neutral/averse. The sample RNG is derived from
$\mathrm{seed}\oplus \mathtt{0xA5A5A5A5}$ and is intentionally
independent of the engine RNG, so editing a cash or inventory value
in the Agents panel and pressing Rebuild produces the same per-tick
trading sequence as a matched run without the edit.

\subsection{Reproducibility}

Every run is driven by a single scalar seed. A mulberry32 generator
$s_{k+1}=(s_k+\mathtt{0x6D2B79F5})\bmod 2^{32}$ produces the
$(0,1)$-uniform stream consumed by every stochastic step---dividend
realisations, zero-intelligence quotes, the noise term in the
utility prior, and the trust-update ordering. Calls to
\texttt{Math.random} from agent or market code are a contract
violation that breaks reproducibility by the (population, seed)
pair and are flagged as an invariant in the project style guide.
The seed is re-rolled only on Reset, and the replay system
reconstructs past states by slicing append-only history arrays on
\texttt{Market} and \texttt{Logger}---no frame is ever recomputed.
A single seed therefore reproduces the entire 720-tick session
end-to-end, rounds and round-end resets included.

Plans~II and~III add a single exogenous source of randomness---the
language-model reply---which the engine handles by keeping the
Plan~I fallback closed-form (\S\ref{sec:plan2},~\S\ref{sec:plan3}).
Paired comparisons across plans therefore always carry a
(plan,~seed) tag, and a Plan~II run with zero successful LLM calls
is definitionally identical to the matching Plan~I run.
